\section{Technical Preparation}

\subsection{Resolution of singularities}

    \begin{theorem}\label{thm:log_resolution_in_characteristic_zero_for_varieties}
        Let \(X\) be a normal variety over an algebraically closed field \(\kkk\) of characteristic zero.
        Then there exists a log resolution of singularities \(f: Y \to X\) such that \(Y\) is smooth and the exceptional divisor \(E\) is a simple normal crossing divisor.
    \end{theorem}

\subsection{Negativity Lemma}

    \begin{theorem}\label{thm:negativity_lemma}
        Let \(f: Y \to X\) be a proper birational morphism between normal varieties.
        Let \(D\) be a \(\bbQ\)-Cartier \(\bbQ\)-divisor on \(Y\) such that \(-D\) is \(f\)-nef.
        Then \(D\) is effective if and only if \(f_*D\) is.
    \end{theorem}
    \begin{proof}
        \Yang{To be completed.}
    \end{proof}

\subsection{General adjunction formula}

    \begin{theorem}[Adjuction formula]\label{thm:general_adjuction_formula}
        Let \(X\) be a normal variety and \(S\) be a reduced divisor on \(X\).
        % Let \(f: Y \to X\) be a log resolution of \((X,S)\).
        % Write \(K_Y + S' = f^*(K_X + S) + E\), where \(S'\) is the strict transform of \(S\) and \(E\) is exceptional.
        % Then \(f_* \calO_Y(\lceil E \rceil) \cong \calO_X\).
        \Yang{Need to check the statement.}
    \end{theorem}
    \begin{proof}
        \Yang{To be completed.}
    \end{proof}

\subsection{Exceptional divisors}

    \begin{proposition}\label{prop:pushout_of_O_X_E_with_E_effective_exceptional_along_birational_morphism}
        Let \(f: Y \to X\) be a proper birational morphism between normal varieties.
        Let \(E\) be an effective \(f\)-exceptional divisor on \(Y\).
        Then we have \(f_* \calO_Y(E) \cong \calO_X\).
    \end{proposition}


\subsection{Base locus}

    \begin{example}
        Let $E$ be an elliptic curve and $P \in E(\kkk)$.
        Set $X = \bbP(\calO_E \oplus \calO_E(-P))$ with the natural projection $\pi: X \to E$ and the unique minimal section $C_0$ with $C_0^2 = -1$.
        Let $D = C_0 + \pi^*P$.
        Then 
        \[ \Bs |D| = \{Q\}, \]
        where $Q$ is some point on $X$ projecting to $P$.

        Indeed, we have 
        \[ h^0(X, D) = h^0(E, \pi_*\left(\calO_X(C_0) \otimes \pi^*\calO_E(P) \right)) = h^0(E, \calO_E \oplus \calO_E(P)) = 2. \]
        Thus $\dim |D| = 1$.
        For $D_t \in |D|$, write $D_t = H_t + \pi^*B$ where $H_t$ is the horizontal part and $B \in \Div(E)$ is the vertical part.
        Intersecting with fibers, we get that $H_t$ is prime and a section of $\pi$.
        By Bernstein's theorem, for general $t$, $D_t$ is irreducible and thus $D_t = H_t$.
        Since $D^2 = 1$, we only need to find the intersection of two distinct members of $|D|$ to find the base locus.
        Note on $\CH_0(X)$, we have 
        \[ D_t \cap D_s = (C_0 + \pi^*P) \cdot (C_0 + \pi^*P) = [Q] \]
        for a point $Q$.
        By proper pushforward of $\CH_0$, we see that $Q$ projects to $P$.
        It follows that any two irreducible members of $|D|$ intersect at $Q$ since they are section and $\Bs |D| \subset \{Q\}$.
        
        Fix an irreducible member $D_0 \in |D|$.
        For every $D_t = H_t + \pi^*B_t \in |D|$, $H_t$ is different from $D_0$ and thus 
        \[ D_t \cap D_0 = [H_t \cap D_0] + [\pi^*B_t \cap D_0] = [Q]. \]
        Note both $[H_t \cap D_0]$ and $[\pi^*B_t \cap D_0]$ are effective, we get that $[\pi^*B_t \cap D_0] = [Q]$ or $0$.
        If $[\pi^*B_t \cap D_0] = [Q]$, then $B_t = P$ and thus $D_t = H_t + \pi^*P \sim C_0 + \pi^*P$.
        Then $H_t \sim C_0$ and thus $H_t = C_0$.
        It pass through $Q$.
        Otherwise, $[\pi^*B_t \cap D_0] = 0$ and thus $B_t = 0$ and thus $D_t$ is irreducible.
        We have seen that $D_t$ also pass through $Q$.
        Hence, $\Bs |D| = \{Q\}$.

        % Hence $Q$ is the unique base point of $|D|$.
    \end{example}
    
